\chap{Certifying an inexact protected correction}\label{sec:inexact-repair}
An existence theorem for a regularized optimizer does not guarantee that a numerical procedure has reached it. The procedure may stop early, use a poorly scaled basis, or report the residual of the wrong equation. These are mathematical distinctions, not implementation details that can be deferred until after a certificate has been interpreted.

The relevant question is now more specific. Given a proposed correction, what quantities can be checked without first computing the exact optimizer? The answer is a normal-equation residual and, independently, a primal and dual objective gap. Both give quantitative statements about the declared regularized problem. Neither promises that the original unregularized realization problem is solvable.

\begin{scopebox}
The current formal extension constructs the positive-parameter optimizer rather than assuming it exists. It proves the energy identity, global uniqueness, and the estimates in this chapter, with a separate bridge to the actual protected kernel. The weak compactness, infinite-limit, and diagonal asymptotic arguments retain their separately stated verification scopes. The exact Python certificate verifier is tested, not universally verified as a refinement of Lean.
\end{scopebox}

\subsection{There are three different residuals}
Keep $H=\ker p$, $T=d|_H$, and a residual datum $r$. Fix $\alpha>0$. Write $a_\alpha$ for WR06's unique minimizer and let $a\in H$ be any proposed correction. Three quantities now have different meanings:
\[
 v=r-Ta,\qquad
 g=T^*(Ta-r)+\alpha a,\qquad
 e=a-a_\alpha.
\]
The vector $v$ is the fit residual. The vector $g$ is the normal-equation residual, equal to one half of the derivative of the quadratic objective after identifying the dual with $H$. The vector $e$ is the unknown correction error. The exact optimizer generally has $v\ne0$ but always has $g=0$.

The adjoint in $g$ is the restricted adjoint on $H$, not the unrestricted ambient adjoint. Equivalently,
\[
 T^*v=P_Hd^*v.
\]
Thus the correction certificate must measure the part of the normal equation that belongs to the permitted correction space. A numerical program that forgets this projection can incorrectly reject an exact protected optimizer.

\subsection{A posteriori correction and fit bounds}
\begin{theorem}[Normal residual as a solver certificate]\label{WR13}
For every proposed $a\in H$,
\begin{align}
 \norm{a-a_\alpha}&\le\frac{\norm g}{\alpha},\label{eq:solver-error}\\
 \norm{T(a-a_\alpha)}&\le\frac{\norm g}{2\sqrt\alpha},\label{eq:solver-fit}\\
 0\le J_\alpha(a)-J_\alpha(a_\alpha)&\le\frac{\norm g^2}{\alpha}.
 \label{eq:solver-objective}
\end{align}
\end{theorem}
\begin{proof}
Put $B=T^*T+\alpha I$. The exact normal equation gives $Be=g$, hence
\begin{equation}\label{eq:solver-energy}
 \norm{Te}^2+\alpha\norm e^2=\langle g,e\rangle\le\norm g\norm e.
\end{equation}
Dropping the first term and dividing when $e\ne0$ proves \eqref{eq:solver-error}. The zero case is immediate. Completing the square in \eqref{eq:solver-energy} gives
\[
 \alpha\norm{Te}^2
 \le \alpha\norm g\norm e-\alpha^2\norm e^2
 \le\frac{\norm g^2}{4},
\]
which proves \eqref{eq:solver-fit}. WR06 and \eqref{eq:solver-energy} identify the objective gap with $\langle g,e\rangle$. Cauchy--Schwarz and \eqref{eq:solver-error} prove \eqref{eq:solver-objective}.

\end{proof}
\written{WR13}

\paragraph{Sharp scalar examples.} For the first sharpness statement take $T=0$ on $\R$ and a nonzero error. Then $g=\alpha e$. For the second, take $T=\sqrt\alpha\,I$ on $\R$, so $g=2\alpha e$ and $\norm{Te}=\norm g/(2\sqrt\alpha)$. These examples establish worst-case sharpness, not simultaneous equality for every operator.

For a fixed raw vector $f$, the corresponding repaired candidates are $f-a$ and $f-a_\alpha$. Their distance is exactly $\norm e$. Both preserve $pf$ because both corrections belong to $H$. The theorem therefore controls a genuine protected candidate error, not just an intermediate coordinate vector.

\fixedplate{37_solver_amplification}{A fixed normal-equation tolerance has different consequences at different regularization parameters. The curves are universal certified upper bounds for a declared normal-residual norm, not observed errors of a particular experiment. The correction estimate grows as $\alpha^{-1}$ while the fitted-data estimate grows as $\alpha^{-1/2}$.}{fig:solver-amplification}

\subsection{Why a vanishing normal residual is not enough}
Let $T=0$ on $\R$, let $r=0$, and take $\alpha_n=n^{-2}$. The exact optimizer is zero for every $n$. The fixed proposal $a_n=1$ instead has
\[
 v_n=0,\qquad g_n=\alpha_n\longrightarrow0,\qquad
 J_{\alpha_n}(a_n)=\alpha_n\longrightarrow0,
 \qquad\abs{a_n-a_{\alpha_n}}=1.
\]
Thus the fit residual, the normal residual, and even the absolute objective value can all tend to zero without convergence to the selected optimizer. The certificate correctly reports $\abs{g_n}/\alpha_n=1$.

This is not an infinite-dimensional obstruction. It occurs in one dimension because the regularizing penalty itself disappears. The question is not whether the reported residual is small in isolation, but whether it is small relative to the parameter that converts it into an error bound.

\subsection{A dual certificate does not require solving the normal equation}
For an arbitrary $z\in Y$, define
\[
 \mathcal D_\alpha(z)=2\langle r,z\rangle-\norm z^2
                    -\frac{\norm{T^*z}^2}{\alpha}.
\]
Unlike the fixed-budget witness in WR04, this dual witness has no unit-norm constraint. Its quadratic penalties supply the control.

\begin{theorem}[Exact regularized primal and dual gap]\label{WR14}
For every $a\in H$ and $z\in Y$,
\begin{equation}\label{eq:regularized-duality-gap}
 J_\alpha(a)-\mathcal D_\alpha(z)
 =\norm{Ta-r+z}^2+
   \frac{\norm{T^*z-\alpha a}^2}{\alpha}.
\end{equation}
Consequently $\mathcal D_\alpha(z)\le J_\alpha(a)$ for every pair. At
$z_\alpha=r-Ta_\alpha$, equality holds, so the dual maximum is attained and equals the primal minimum. Moreover,
\[
 \alpha\norm{a-a_\alpha}^2
 \le J_\alpha(a)-\mathcal D_\alpha(z).
\]
For the easily computed choice $z=r-Ta$,
\begin{equation}\label{eq:residual-dual-gap}
 J_\alpha(a)-\mathcal D_\alpha(r-Ta)=\frac{\norm g^2}{\alpha}.
\end{equation}
\end{theorem}
\begin{proof}
Expand the two squares in \eqref{eq:regularized-duality-gap}. The mixed terms containing $Ta$ and $z$ cancel by the adjoint identity. The remaining terms are exactly the left side. Nonnegativity gives weak duality. At the exact optimizer, the first square vanishes for $z_\alpha$, and the second vanishes by the normal equation. This supplies a dual optimizer directly, without exchanging an infimum and a supremum.

Weak duality at $a_\alpha$ gives
\[
 J_\alpha(a)-\mathcal D_\alpha(z)
 \ge J_\alpha(a)-J_\alpha(a_\alpha)
 \ge\alpha\norm{a-a_\alpha}^2
\]
by WR06. Finally $z=r-Ta$ makes the first square zero and the second numerator equal to $\norm{-g}^2$, proving \eqref{eq:residual-dual-gap}.
\end{proof}
\written{WR14}

A proposed correction and a proposed dual witness can come from different procedures. The verifier only needs the mathematical problem, the two vectors, and the displayed identities. A poor witness remains valid as a lower bound, but can give a loose gap. The residual witness avoids a second optimization routine and recovers WR13's bound exactly.

\subsection{A protected example with every certificate calculated}
Take
\[
 d(x,y)=x,\qquad p(x,y)=x-y,\qquad f=(1,0),\qquad r=df=1,\qquad\alpha=1.
\]
Every allowed correction has the form $(t,t)$. The objective is
\[
 J(t,t)=(t-1)^2+2t^2,
\]
so $a_\alpha=(1/3,1/3)$ and $J(a_\alpha)=2/3$. Suppose a procedure returns only $a=(1/4,1/4)$. Then
\[
 v=\frac34,\qquad T^*v=(3/8,3/8),\qquad
 g=(-1/8,-1/8),\qquad\norm g^2=\frac1{32}.
\]
The actual correction-error square is $1/72$, while the certified upper bound is $1/32$. The actual fitted-error square is $1/144$, while the sharper bound in WR13 is $1/128$.

The proposed objective is $11/16$. Its actual excess over the optimum is $1/48$. The residual witness gives
\[
 \mathcal D_1(3/4)=\frac{21}{32},\qquad
 J(a)-\mathcal D_1(3/4)=\frac1{32}.
\]
The duality gap is an upper bound for suboptimality, not necessarily the exact suboptimality. Both repaired candidates preserve $p(f)=1$.

\fixedplate{38_objective_certificate}{The objective on the actual protected line is $3t^2-2t+1$. The proposed correction at $t=1/4$, exact minimizer at $t=1/3$, and residual-witness lower bound are distinct. The vertical gap to the dual level is the checkable certificate. The smaller gap to the minimum is the true suboptimality.}{fig:objective-certificate}

The supplied program \path{examples/inexact_repair_certificate.py} verifies this calculation without calling the optimizer. It also accepts a user-supplied dual witness. Every reported scalar is rational, and all norm error bounds are stored as squared quantities so that no rounding of a square root is required.
